\begin{figure}[ht]
   \centering
   \begin{tikzpicture}[
         nodestyle/.style={
               fill=black,
               circle,
               minimum height=4pt,
               inner sep=0pt,
            }
      ]
      \draw[->](-4,0) -- (4,0) node[right]{$F1$};%x軸
      \draw[->](0,-4) -- (0,4) node[above]{$F2$};%y軸
      \node[nodestyle] (item1) at (1.5,3.1) [] {};
      \node[nodestyle] (item2) at (2.1,2.8) [] {};
      \node[nodestyle] (item3) at (2.2,2.1) [label=item3] {};
      \node[nodestyle] (item4) at (1.5,-2.0) [label=item4] {};
      \node[nodestyle] (item5) at (1.5,-2.5) [label=item5] {};
      \node[nodestyle] (item6) at (-1.5,-2.8) [label=item6] {};
      \draw[->,thick](-2,-3.464) -- (2,3.464) node[above right]{$F1^\prime$};%回転後
      \draw[->,thick](3.464,-2) -- (-3.464,2) node[above left]{$F2^\prime$};%回転後
      \draw[dashed](0,0) -- (3/2,3*1.732/2);
      \draw [->](1.5,0) arc [start angle = 0, end angle = 60, radius = 1.5];
      \draw [->](0,1.5) arc [start angle = 90, end angle = 150, radius = 1.5];
      \coordinate (o) at (0,0);
      \coordinate (F1x) at (2,3.464);
      \coordinate (F2x) at (-3.464,2);
      \draw[dashed] (item1) -- ($(o)!(item1)!(F1x)$){};
      \draw[dashed] (item1) -- ($(o)!(item1)!(F2x)$){};
   \end{tikzpicture}
   \caption{因子軸の回転。60度回転させてみました。}
   \label{fig::x10_FAcoord2}
\end{figure}
